\documentclass[border=9pt]{standalone}
\usepackage[fontset=fandol]{ctex}
\usepackage{amsmath,amssymb,tikz}
\usetikzlibrary{arrows.meta,calc,decorations.pathmorphing}
\tikzset{arr/.style={-{Stealth[length=2mm]},semithick}}
\begin{document}

\begin{tikzpicture}[scale=.9]
\draw[thick] (-4.5,0)--(4.7,0) node[right] {地面};
\draw[thick] (-4.5,2)--(4.7,2) node[right] {电离层底面};
\draw[blue] (2,0) arc(0:180:2);
\draw[red,thick] (3.2,0) arc(0:180:3.2);
\draw[gray,dashed] (4,0) arc(0:180:4);
\fill (0,0) circle(2pt) node[below] {$O$};
\fill (0,2) circle(2pt) node[above left] {$A$};
\fill (2.498,2) circle(2pt) node[above right] {$B$};
\draw[arr,red] (0,0)--(2.498,2) node[midway,below right] {$ct$};
\draw[<->] (-4.1,0)--(-4.1,2) node[midway,left] {$H_i$};
\draw[<->] (0,1.55)--(2.498,1.55) node[midway,below] {$r_i$};
\draw[arr] (2.55,2.1)--(3.46,2.1);
\node at(.1,4.65) {球面波以$c$传播，交线的径向速度可大于$c$};
\node at(.1,-.8) {$r_i^2+H_i^2=c^2t^2$};
\end{tikzpicture}

\end{document}
